\documentclass[11pt,a4paper]{article}
% ===================================================================
%  Vorlage: Arbeitsblatt „Matt in 1“ mit Lösungsteil
%  Jede Stellung vorher prüfen:  python3 pruefe_stellung.py "FEN"
%  Kompilieren: pdflatex vorlage_matt_in_eins.tex
% ===================================================================
\usepackage[T1]{fontenc}
\usepackage{lmodern}
\usepackage[ngerman]{babel}
\usepackage[a4paper,margin=2cm]{geometry}
\usepackage{xcolor}
\usepackage{xskak}
\usepackage{chessboard}
\pagestyle{empty}
\setlength{\parindent}{0pt}

% Einheitliches Aussehen aller Bretter dieses Blatts
\setchessboard{smallboard, showmover=true, label=true}

% \Aufgabe{Nummer}{FEN}{Hinweis}
\newcommand\Aufgabe[3]{%
  \begin{minipage}[t]{0.48\textwidth}\centering
    \textbf{Aufgabe #1}\par\smallskip
    \chessboard[setfen={#2}]\par\smallskip
    {\small #3}\par\smallskip
    Lösung: \rule{3cm}{0.4pt}
  \end{minipage}}

\begin{document}
{\Large\bfseries\sffamily Matt in einem Zug} \hfill Name: \rule{5cm}{0.4pt}

\smallskip
Weiß ist am Zug und setzt mit \emph{einem} Zug matt. Schreibe den Zug auf.
Das kleine Quadrat neben dem Brett zeigt, wer am Zug ist.

\bigskip
\Aufgabe{1}{6k1/5ppp/8/8/8/8/5PPP/3R2K1 w - - 0 1}{Die eigenen Bauern sind im Weg.}\hfill
\Aufgabe{2}{r1bqkbnr/pppp1ppp/2n5/4p3/2B1P3/5Q2/PPPP1PPP/RNB1K1NR w KQkq - 4 4}{Ein schwacher Punkt: f7.}

\bigskip\bigskip
\Aufgabe{3}{6rk/6pp/8/6N1/8/8/8/6K1 w - - 0 1}{Der König erstickt.}\hfill
\Aufgabe{4}{1k6/ppp5/8/8/8/8/8/K6R w - - 0 1}{Die letzte Reihe.}

\vfill
\rule{\textwidth}{0.4pt}\par
{\small\textbf{Lösungen:}
1)~\symrook d8\# \quad 2)~\symqueen xf7\# \quad
3)~\symknight f7\# \quad 4)~\symrook h8\#}
\end{document}
